\section{Les Bases Python}

\subsubsection{Opérateurs et Affectation}

\begin{itemize}
\item Affectation \texttt{=} : Assigne une valeur (à droite) à une variable (à gauche).
\item Addition \texttt{+} : Somme
\item Soustraction \texttt{-} : Différence
\item Multiplication \texttt{*} : Produit.
\item Division \texttt{/} : Division classique (renvoie toujours un `float`).
\item Division entière \texttt{//} : Renvoie uniquement la partie entière du quotient.
\item Modulo \texttt{\%} : Renvoie le reste de la division entière.
\item Puissance \texttt{**} : Élève à l'exposant.
\end{itemize}

\subsubsection{Affectation Combinée}
Modifie une variable en utilisant sa propre valeur. On a \texttt{+=}, \texttt{-=}, \texttt{*=}, \texttt{/=}, $\cdots$ . C'est équivalent à \texttt{x = x + y}.

\subsubsection{Opérateurs de Comparaison}
Renvoient \texttt{True} ou \texttt{False} :

\begin{itemize}
\item \texttt{==} : Égal à
\item \texttt{!=} : Différent de.
\item \texttt{<} : Strictement inférieur à.
\item \texttt{>} : Strictement supérieur à.
\item \texttt{<=} : Inférieur ou égal à.
\item \texttt{>=} : Supérieur ou égal à.
\end{itemize}

\subsubsection{Types de Données Fondamentaux}

On peut voir le type de une variable avec la commande \texttt{type(variable)}

\begin{itemize}
\item Nombre entier \texttt{int} : Chiffre entier
\item Nombre décimal \texttt{float} : Utilisation du point
\item Valeur de vérité logique \texttt{bool} : \texttt{True} ou \texttt{False}
\item Chaîne de caractères \texttt{str} : Délimité par \texttt{"..."} ou \texttt{'...'}
\item Liste \texttt{list} : Délimité par \texttt{[...]}
\item Dictionnaire \texttt{dict} : Délimité par \texttt{\{...\}}
\end{itemize}

\subsubsection{Opérateurs Logiques}

\begin{itemize}
\item \texttt{and} : Toutes les conditions doivent être vraies.
\item \texttt{or} : Au moins une des conditions doit être vraie.
\item \texttt{not} : Inverse le résultat logique
\end{itemize}